\paragraph{How are prescription approval decisions distributed across the study populations?}
Prescription decisions for the study populations.

\begin{table}[htbp]
\centering
\caption{Prescription decisions for the study populations.}
\label{tab:study-population-prescription-decisions}
\small
\begin{tabular}{@{}>{\raggedright\arraybackslash}p{0.45\linewidth}rrrr@{}}
\toprule
Population & Prescriptions & Approved & Unknown & Rejected \\
\midrule

All approval channels, matched to first-order lines & 103,550 & 36,908 & 66,595 & 47 \\

\texttt{CUSTOMER\_UPLOAD}, matched to first-order lines & 95,106 & 30,537 & 64,569 & 0 \\

\texttt{CUSTOMER\_UPLOAD} aged-cancellation subset & 26,443 & 6 & 26,437 & 0 \\
\bottomrule
\end{tabular}

\begin{minipage}{\linewidth}
\footnotesize
\textit{Notes:} Distinct prescriptions in the same populations as
Table~\ref{tab:study-population-counts}. Approved means
\texttt{PRESCRIPTION\_APPROVED = TRUE}, consistent with the prescription registry's
\texttt{APPROVED\_NUMERATOR}. Rejected means current \texttt{RX\_STATUS = DENIED}
without an approval flag. Unknown includes cancelled and pending prescriptions
without a recorded approval or rejection outcome. The three categories are
mutually exclusive and sum to the prescription count in each row.
Decisions use the current source snapshot refreshed November 04, 2026;
approval and denial times are not restricted to the aged-cancellation cutoff.
\end{minipage}
\end{table}
